\subsubsection{Central coordinates, orientation, and regional records}
\label{el-centro:registros}

The central structure relates position, evaluation, and transport. At
cell \((5,5)\), the two APP sheets give
\begin{equation}
 \mathcal A(5,5)=\bigl((1,1),(7,2)\bigr),\qquad
 5+5=1+9\cdot1,\qquad5\cdot5=7+9\cdot2.
 \label{el-centro:app}
\end{equation}
The two ones on the additive sheet are, in order, residue and quotient.
Their numerical coincidence preserves two distinct arithmetic functions.
On the sum-\(10\) fiber, the two oriented vacancies satisfy
\begin{equation}
 \mathcal A(8,2)=\mathcal A(2,8)=\bigl((1,1),(7,1)\bigr).
 \label{el-centro:vacancias}
\end{equation}
The complete additive sheet and the multiplicative residue are preserved;
the quotient coordinate of the product decreases by exactly
\(1\). Transport additionally retains which of the two conjugate
positions has been reached.

\begin{figure}[H]
\centering
\begin{tikzpicture}[x=.67cm,y=.67cm,font=\footnotesize,>=Latex]
 \draw[step=1,black!12,line width=.3pt] (1,1) grid (9,9);
 \draw[->,black!60] (.7,1)--(9.5,1) node[right] {$i$};
 \draw[->,black!60] (1,.7)--(1,9.5) node[above] {$j$};
 \foreach \k in {1,...,9}{
  \node[below=2pt] at (\k,1) {$\k$};
  \node[left=2pt] at (1,\k) {$\k$};
 }
 \draw[black!65] (1,9)--(9,1);
 \draw[<->,line width=.9pt] (2,8)--(8,2);
 \foreach \x/\y in {2/8,5/5,8/2}{
  \filldraw[fill=white,line width=.8pt] (\x,\y) circle (3pt);
 }
 \fill (5,5) circle (1.5pt);
 \node[above right=3pt,fill=white,inner sep=1pt] at (2,8) {$(2,8)$};
 \node[above right=3pt,fill=white,inner sep=1pt] at (5,5) {$(5,5)$};
 \node[above right=3pt,fill=white,inner sep=1pt] at (8,2) {$(8,2)$};
 \node[anchor=west,align=left] at (10.5,7.6)
   {$i+j=10$\\$R^+=1,\quad q^+=1$};
 \node[anchor=west,align=left] at (10.5,5.4)
   {center: $ij=25$\\$R^\times=7,\quad q^\times=2$};
 \node[anchor=west,align=left] at (10.5,3.1)
   {vacancies: $ij=16$\\$R^\times=7,\quad q^\times=1$};
\end{tikzpicture}
\caption{Electronic center and conjugate vacancies in the APP chart.
The diagonal is the additive fiber \(\mathcal F_{10}\); the center
is the fixed point of \((i,j)\mapsto(10-i,10-j)\). The two
shifts of magnitude \(3\) preserve the sum and product
residue and reduce the multiplicative quotient by one unit. The
arrows represent deformations on the fiber.}
\label{el-centro:figura}
\end{figure}


Central symmetry also characterizes its persistence under changes of
chart. If an admissible atlas transformation \(f\) intertwines
inversion, \(f\rho=\rho f\), then
\[
 \rho f(5,5)=f\rho(5,5)=f(5,5).
\]
Uniqueness of the fixed point forces \(f(5,5)=(5,5)\). This is the
internal property making the central location stable: every
transformation of that type preserves the cell before evaluating its
electronic character.

The central matrix block belongs to evaluation on the two
adjacent positions \(4,5\). Its decomposition, established in the
local development of the kernel, is
\[
 B_c=\begin{pmatrix}7&2\\2&7\end{pmatrix}=9P_++5P_-.
\]
Here \(9\) and \(5\) are eigenvalues of two modes, whereas
\((5,5)\) is a position in the atlas. Concatenations of its
entries give \(72+27=77+22=99\) and
\(77-22=55=54+1\). The first equality preserves the total sum under
two arrangements; the second expresses the defect between diagonal
and antidiagonal readings. These operations precede electronic
composition and preserve their matrix provenance.

The initial orientation is another datum. In the representation of
electronic return, \((h_0,v_0)=(1,1)\) is fixed, that is, positive
advance in both coordinates. Simultaneous inversion every \(27\)
steps, read in windows of \(9\), determines
\[
 \tau_e=(+1,+1,+1,-1,-1,-1,+1,+1,+1,-1,-1,-1).
\]
The segments \((1,1,1)\) belong to this orientation sequence:
each segment comprises three windows in the same direction. The recurrence
in the section on readers proves this record in
\eqref{eq:el-tau}, and its two evaluations recover the triple
\((80,54,6)\) entering
\eqref{eq:el-formula-completa}. The positional coordinate,
APP quotients, and sign sequence are thus connected through
explicit operations.

The signature \(555555\) belongs to a regional reading in another
domain. For a decimal emission \(U=(U_1,\ldots,U_6)\), we write
\(U_j=100d_{1j}+10d_{2j}+d_{3j}\), with three digits, including
leading zeros. The multiplicative-signature reader is
\[
 \mathcal E_\times(U)
   =\bigl([d_{2j}-d_{1j}]_{10}\bigr)_{j=1}^{6}.
\]
In the transversal region
\(U_\pi^\perp=(501,614,498,169,272,272)\), the differences are
\((-5,-5,5,5,5,5)\). Consequently,
\begin{equation}
 \mathcal E_\times(U_\pi^\perp)=(5,5,5,5,5,5),\qquad
 \Psi(U_\pi^\perp)=010211.
 \label{el-centro:firma-transversal}
\end{equation}
The first map preserves a multiplicative signature of the catalog;
the second preserves the ternary word common to the five regions.
Figure \ref{pi-estrella:figura} shows the transversal position of
this emission and its multiplicity \(432\).

\paragraph{Memory required to transport the constant signature.}
The signature \(555555\) additionally provides a verifiable example of
the trajectory information required by TPK. Let
\(\mathcal X_\times=D_9^2\times\{N,E,S,O\}\), with \(324\)
oriented positions in the product channel. Advance \(P\) moves
one square in the stored direction, and the half-turn \(H\)
reverses that direction. In this representation, the reader
\(E_{\rm sig}\) sums the discrete exponents by windows:
\[
 \ell_9(1,2,4,8,7,5)=(0,1,2,3,4,5),\qquad
 \ell_9(3)=\ell_9(6)=\ell_9(9)=0.
\]
During each of the six nine-step windows,
\(\ell_9(\rho_9(ij))\) is evaluated before every product advance, selected
by TRIT at times \(2,5,8\pmod9\). The window sum is
reduced modulo \(10\). After steps \(27\) and \(54\), the
direction is reversed. This rule determines the six components of
\(E_{\rm sig}\) from the initial position and orientation.

Enumeration of this domain produces \(26\) signatures. The fiber of
\(555555\) contains \(24\) oriented positions; after one
advance, their images are distributed among three signatures, with eight
representatives each:
\[
 555177,\qquad555717,\qquad555771.
\]
Three witnesses, in coordinates from \(1\) to \(9\), allow the
distinction to be checked directly:
\[
\begin{array}{c|c|c}
x&E_{\rm sig}(x)&E_{\rm sig}(Px)\\\hline
(3,5,N)&555555&555177\\
(3,4,N)&555555&555717\\
(3,4,S)&555555&555771
\end{array}
\]
Thus, the common value of the signature does not determine its next value:
the transport class within that fiber forms part of the operational
datum. The minimal stable partition is calculated by starting from
signature equivalence \(\sim_0\) and refining through
\[
 x\sim_{n+1}y\quad\Longleftrightarrow\quad
 x\sim_n y,\quad Px\sim_n Py,\quad Hx\sim_n Hy.
\]
The procedure terminates because the domain is finite. Every stable
congruence refining \(\sim_0\) refines every \(\sim_n\), by induction;
its fixed point is therefore the minimal stable congruence required.
The census gives \(26\to28\to28\), with histogram
\(27\cdot8+108=324\). The only signature that splits is
\(555555\), into three classes of eight elements. The accompanying certificate
reproduces the enumeration, the three classes, and stability under both
operators.

The multiplicity \(24\) here belongs to the product-channel fiber;
the multiplicity \(432\) belongs to the complete regional emission
\(U_\pi^\perp\). The first proves why the signature needs transport
memory; the second places that signature within the joint catalog of
the two sheets. Both readings participate in the same architecture,
with specified domains and operations.

This distinction of domains makes it possible to bring together the occurrences of unity
and the center with their respective meanings:
\begin{center}
\small
\begin{tabularx}{\textwidth}{@{}lX@{}}
\toprule
Record & Object and function \\
\midrule
\((5,5)\) & Fixed cell of the central inversion of the APP atlas.\\
\((R^+,q^+)=(1,1)\) & Residue and quotient of the central additive evaluation.\\
\((h_0,v_0)=(1,1)\) & Initial orientation of the two shifts.\\
\((+1,+1,+1)\) & Three consecutive windows of \(\tau_e\) with positive orientation.\\
\(c=111111\) & Datum of the fifth nonadic boundary, Section \ref{mono-recta:desarrollo}.\\
\(\mathcal E_\times=555555\) & Signature of the transversal region of \(\pi\), with six positions.\\
\bottomrule
\end{tabularx}
\end{center}

Electronic coordination is obtained by composing these levels in their
dependency order. APP determines the center and evaluations;
TRIT selects regimes and orientations; TPK prolongs the trajectory
and preserves its records. The previously generated regions provide
the characters \(\pi,\varphi,e\); the dodecaphase record provides
\(\alpha\); the electronic readers determine the factors of
equation \eqref{eq:el-formula-completa}. This composition preserves
the identity of the section throughout its different
representations, up to dimensional evaluation and subsequent
metrological comparison.
